\title{FlashAttention-3: Fast and Accurate Attention with Asynchrony and Low-precision}

\begin{abstract}
Attention mechanisms serve as the foundational component of the widely adopted Transformer architecture, yet they remain the primary computational bottleneck in scaling large language models and accommodating longer input contexts. \textsc{FlashAttention} introduced a strategy to accelerate attention computations on GPUs by significantly reducing memory input/output traffic. Nevertheless, previous implementations have not fully leveraged recent advances in hardware; for instance, \textsc{FlashAttention-2} utilizes only 35\% of the H100 GPU’s computational capacity. To address this, we present three key optimizations for attention on Hopper GPUs: (1) utilizing the asynchronous capabilities of Tensor Cores and TMA through warp specialization to concurrently perform computation and data transfer, (2) fusing block-wise matrix multiplication with softmax computation via interleaving, and (3) employing block quantization alongside incoherent execution that taps into hardware-supported FP8 low-precision arithmetic. Our method, \textsc{FlashAttention-3}, achieves a 1.5-2.0$\times$ performance increase on H100 GPUs; FP16 execution attains up to 740 TFLOPs/s (reflecting 75\% utilization), while FP8 mode approaches 1.2 PFLOPs/s. We further show that FP8 \textsc{FlashAttention-3} yields 2.6$\times$ smaller numerical errors relative to a standard FP8 attention implementation.
\end{abstract}

\section{Introduction}
\label{sec:intro}

Within the Transformer framework~\citep{vaswani2017attention}, the principal source of computational overhead lies in the attention mechanism, as the process of calculating self-attention scores between queries and keys requires operations that grow quadratically with respect to sequence length. Extending the attention mechanism to support longer contexts has the potential to enable a variety of novel capabilities, such as reasoning across multiple extended documents~\citep{guo2021longt5,shaham2022scrolls,peng2023yarn}, managing extensive files in sizeable code repositories~\citep{roziere2023code, li2023starcoder}, addressing new input modalities like detailed images~\citep{chen2022scaling}, audio signals~\citep{gulati2020conformer}, and video data~\citep{ho2022video}, as well as supporting new use cases involving sustained user interactions~\citep{sun2019bert4rec} and complex, long-term agent tasks~\citep{yao2022react}. These opportunities have sparked a considerable focus on accelerating attention computation for long sequences. Efforts in this direction include developing approximation methods~\citep{katharopoulos2020transformers,choromanski2020rethinking, tay2020efficient}, enhancing software performance (\citep{rabe2021self,dao2022flashattention,kwon2023efficient}), and exploring fundamentally different model architectures~\citep{peng2023rwkv,sun2023retentive,gu2023mamba}.

This study extends upon the foundational work of \citet{dao2022flashattention}, who designed exact-attention algorithms by factoring in the specifics of GPU execution models and hardware traits at the algorithmic level. In their research, Dao et al. presented \textsc{FlashAttention}, a pioneering tiling mechanism to parallelize attention operations, thereby eliminating the need for slow global memory accesses by combining all attention-related steps within a single GPU kernel.

Following this, \citet{dao2023flashattention2} introduced \textsc{FlashAttention-2}, redesigning the algorithm to also parallelize computations across the sequence length and organize the forward-pass inner loop around blocks within the key and value matrices. This modification aimed to enhance both the GPU's workload allocation and overall occupancy. Despite these refinements, our analysis reveals that \textsc{FlashAttention-2} still suffers from suboptimal utilization on the latest GPU architectures, especially when compared to highly optimized matrix-multiplication (GEMM) kernels; for instance, on Hopper H100 GPUs, utilization rates are approximately 35\%, notably less than the 80-90\% observed for GEMM. Some of these discrepancies are likely caused by implementation details, such as the absence of Hopper-specific instructions, relying instead on older Ampere commands when leveraging Tensor Cores. Work like ThunkerKitten~\citep{spector2024thunder} and cuDNN 9~\citep{cudnn9} demonstrates that adopting tile-based abstractions and Hopper-optimized instructions can substantially accelerate attention operations and reduce implementation complexity.

At its core, the algorithm underlying \textsc{FlashAttention-2} is structured around a synchronous execution paradigm and does not explicitly integrate asynchrony or low-precision computation within its framework. Asynchrony emerges from hardware advancements that optimize critical operations in machine learning workloads, such as dedicated units for matrix multiplication (Tensor Cores) or specialized modules for efficient memory access (Tensor Memory Accelerator -- TMA), which operate independently from the primary CUDA cores responsible for handling logical, integer, and floating-point tasks. The adoption of reduced precision formats—including FP8 on Hopper and FP4 on Blackwell, as well as earlier standards like FP16 (introduced in Pascal, 2017) and BF16 (implemented in Ampere, 2020)—has consistently enabled significant gains in computational throughput without increasing power consumption or chip size. A comprehensive overview of Hopper's enhancements in this regard is presented in \cref{subsec:hardware}. The key technical hurdle lies in refactoring \textsc{FlashAttention-2} so it can leverage these hardware features: utilizing asynchrony necessitates the concurrent execution of matrix multiplication and softmax operations even when dependencies exist, while operating at low precision demands careful handling to reduce quantization artifacts, a critical concern especially for atypical features found in LLMs~\citep{dettmers2208llm, sun2024massive}.

To address these challenges, we introduce \textsc{FlashAttention-3}, which incorporates three innovative concepts designed to enhance efficiency on modern GPU platforms:\footnote{Our findings are discussed primarily with respect to NVIDIA's Hopper architecture, but the algorithm remains applicable to any GPU system featuring strong asynchronous operations and support for low-precision computation.}

\begin{enumerate}

\item \textbf{Producer-Consumer asynchrony:} We propose a warp-specialized software pipelining strategy that leverages asynchronous operations between data transfer and Tensor Core computation. By partitioning the producer and consumer roles into distinct warps, this approach enhances the algorithm's effectiveness at masking both memory access and instruction issuance delays.

\item \textbf{Hiding softmax under asynchronous block-wise GEMMs:} The low-throughput non-GEMM tasks within softmax, such as floating-point multiply-add and exponentiation, are overlapped with asynchronous WGMMA GEMM instructions. To facilitate this, the \textsc{FlashAttention-2} algorithm is redesigned to eliminate sequential dependency constraints between the softmax procedure and block-wise GEMMs. For instance, in the two-stage variant of the method, softmax processes one score matrix block while WGMMA asynchronously prepares the subsequent block in parallel.

\item \textbf{Hardware-accelerated low-precision GEMM:} The forward pass algorithm is adapted to utilize FP8 Tensor Cores for GEMM operations, achieving an approximate twofold increase in sustained TFLOPs/s. This necessitates resolving the layout compatibility requirements imposed by WGMMA for arranging FP32 accumulator blocks and FP8 operand blocks in memory. To address the precision degradation inherent in FP8 computations, we apply block quantization and incoherent processing techniques.
\end{enumerate}

In order to empirically assess the effectiveness of our approach, we conduct a performance evaluation of \textsc{FlashAttention-3} on the H100 SXM5 GPU across various parameter configurations. Our results indicate that (1) the FP16 variant offers a speed increase of 1.5-2.0$\times$ compared to \textsc{FlashAttention-2} during the forward computation (achieving up to 740 TFLOPs/s), and a 1.5-1.75$\times$ acceleration in the backward computation; (2) FP8 mode attains performance approaching 1.2 PFLOPs/s; and (3) at higher sequence lengths, FP16 consistently demonstrates superior throughput, while FP8 remains a competitive alternative\footnote{Specifically, with a head dimension of 64, FP8 \textsc{FlashAttention-3} surpasses competing methods, but for head dimensions of 128 and 256, it matches their performance without causal masking and falls behind when causal masking is enabled.} compared to NVIDIA's cuDNN attention implementation. We further establish that FP16 \textsc{FlashAttention-3} exhibits identical numerical precision to \textsc{FlashAttention-2} and outperforms the conventional attention implementation, since operations such as softmax rescaling are retained in FP32 precision. Additionally, FP8 \textsc{FlashAttention-3} equipped with block quantization and incoherent processing demonstrates a 2.6$\times$ improvement in accuracy relative to conventional attention with per-tensor quantization, especially in scenarios containing outlier features.

We have released \textsc{FlashAttention-3} under a permissive license\footnote{\textsc{FlashAttention-3}
  can be found at \texttt{https://github.com/Dao-AILab/flash-attention}}, and aim to incorporate it into both PyTorch and Hugging Face libraries so that it may be accessible to a wide community of researchers and developers.

\section{Background: Multi-Head Attention and GPU Characteristics}
\label{sec:background}

\subsection{Multi-Head Attention}
\label{subsec:multi_head_attn}

Let $\mathbf{Q}, \mathbf{K}, \mathbf{V} \in \mathbb{R}^{N \times d}$ be the query, key and value input sequences associated to a single head, where $N$ is the sequence length and $d$ is the head dimension. Then the attention output $\mathbf{O}$ is computed as:

\begin{equation*}
  \mathbf{S} = \alpha \mathbf{Q} \mathbf{K}^\top \in \mathbb{R}^{N \times N}, \quad \mathbf{P} = \mathrm{softmax}(\mathbf{S}) \in \mathbb{R}^{N \times N}, \quad \mathbf{O} = \mathbf{P}\mathbf{V} \in \mathbb{R}^{N \times d},
\end{equation*}

Here, the $\mathrm{softmax}$ function is executed on each row, and commonly, $\alpha$ is chosen as $1/\sqrt{d}$ for the scaling coefficient. To mitigate numerical issues associated with exponentiation, $\mathrm{rowmax}(\mathbf{S})$ is subtracted from $\mathbf{S}$ during the computation. In the context of multi-head attention (MHA), distinct query, key, and value projections are assigned to each head. The calculation can be performed in parallel for several heads and batches, generating the complete output tensor.

Now let $\phi$ be a scalar loss function and let $\mathbf{d}(-) = \partial \phi / \partial (-)$ be notation for the gradient. Given the output gradient $\mathbf{dO} \in \mathbb{R}^{N \times d}$, we compute $\mathbf{dQ}$, $\mathbf{dK}$, and $\mathbf{dV}$ according to the chain rule as follows:

\begin{align*}
  \mathbf{dV} &= \mathbf{P}^\top \mathbf{dO} \in \mathbb{R}^{N \times d} \\
  \mathbf{dP} &= \mathbf{dO} \mathbf{V}^\top \in \mathbb{R}^{N \times N} \\
  \mathbf{dS} &= \mathrm{dsoftmax} (\mathbf{dP}) \in \mathbb{R}^{N \times N} \\
  \mathbf{dQ} &= \alpha \mathbf{dS} \mathbf{K} \in \mathbb{R}^{N \times d} \\
  \mathbf{dK} &= \alpha \mathbf{dS}^\top \mathbf{Q} \in \mathbb{R}^{N \times d},
\end{align*}

In this context, $\mathbf{d}s = (\mathrm{diag}(p) - p p^\top)\mathbf{d}p$ expresses the derivative when $p$ is defined as $\mathrm{softmax}(s)$, considered as a function of the vector $s$. The notation $\mathrm{dsoftmax}(\mathbf{dP})$ refers to utilizing this derivative formula on each row individually. For the backward pass in MHA, this operation remains parallelizable over both the number of heads and batch dimension.

\subsection{GPU hardware characteristics and execution model}
\label{subsec:hardware}

We outline the pertinent features of the GPU execution model that pertain to \textsc{FlashAttention-3}, emphasizing the NVIDIA Hopper architecture as a specific example within this context.

\paragraph{Memory hierarchy:} The arrangement of the GPU's memory follows a layered structure, where memory spaces with smaller capacities typically offer higher bandwidths (\cref{tab:gpu-hierarchy})\footnote{\citet{luo2024benchmarking} provides a shared memory bandwidth estimate of 128 bytes per clock cycle per SM, which we scale by 132 SMs and a boost clock of 1830 MHz for overall bandwidth.}. The largest pool, referred to as global memory (GMEM) or HBM, comprises off-chip DRAM that is universally addressable by all streaming multiprocessors (SMs). Any data accesses from GMEM are automatically served by an on-chip L2 cache. Each individual SM is also equipped with a relatively modest, on-chip, highly banked cache named shared memory (SMEM) that is explicitly managed by the programmer. The closest memory level to compute within an SM is the register file.

\paragraph{Thread hierarchy:} Within the GPU programming paradigm, computational units are systematically structured into threads, which are aggregated at multiple levels. Starting from the most granular, the hierarchy consists of individual threads, which are grouped into warps (each consisting of 32 threads). Four adjacent warps form a warpgroups, and these are further encapsulated into threadblocks (otherwise referred to as cooperative thread arrays or CTAs). In the Hopper architecture, threadblock clusters represent an additional layer of grouping, all of which collectively form a grid at the highest level.

The relationship between these two hierarchies is highly integrated. Threads belonging to the same CTA are scheduled simultaneously on a single SM, while CTAs assigned to the same cluster are simultaneously scheduled on a single GPC. SMEM can be accessed by any thread inside a CTA, whereas individual threads are limited to a maximum of 256 private registers (RMEM) each.

\paragraph{Asynchrony and warp-specialization:}

GPUs are designed as throughput-oriented processors, leveraging high degrees of concurrency and asynchrony to mask both memory access and execution delays. In the context of asynchronous memory transfers between GMEM and SMEM, the Hopper architecture introduces the Tensor Memory Accelerator (TMA), a specialized hardware module \cite[\S7.29]{cuda}. In addition, distinguishing itself from earlier designs such as Ampere, Hopper's Tensor Core—accessed through the warpgroup-wide WGMMA instruction \cite[\S9.7.14]{ptx}—operates asynchronously and has the capability to read its inputs directly from shared memory.

Enabling hardware-driven asynchrony facilitates warp-specialized kernels, where the warps within a CTA are explicitly designated as either producers or consumers, restricting each to solely handle data transfer or computation tasks. This division typically enhances the compiler's capacity to produce more efficient instruction scheduling \citep{warp-specialization-2011}. Furthermore, Hopper includes the ability to dynamically adjust register allocation across warpgroups through the \verb|setmaxnreg| directive \cite[\S9.7.17.1]{ptx}, thereby allowing warps responsible for executing MMAs to utilize a greater portion of RMEM compared to warps engaged in issuing TMA operations, which require only a single thread.

\paragraph{Low-precision number formats:}
\label{sec:low-precision-gpu}
Contemporary GPUs incorporate dedicated hardware for efficient low-precision arithmetic. As an illustration, the WGMMA instruction leverages FP8 Tensor Cores within Hopper architectures, resulting in twice the throughput per SM relative to computations performed with FP16 or BF16 formats.

Successfully utilizing FP8 WGMMA requires a clear understanding of how operand layouts are restricted. When performing GEMM to calculate $A \times B^{\top}$, where $A$ is an $M\times K$ matrix and $B$ is an $N\times K$ matrix, an operand ($A$ or $B$) is referred to as \emph{mn-major} if its elements are contiguous along the outer $M$ or $N$ axis; conversely, it is called \emph{k-major} if the inner $K$ dimension is contiguous in memory. For FP16 WGMMA operations, SMEM operands may be laid out in either mn-major or k-major order. By contrast, FP8 WGMMA restricts operand storage to the k-major format. This restriction introduces challenges—particularly in contexts like attention, where consecutive GEMM operations are fused within one kernel. In these cases, the incompatible layouts between FP32 accumulators and FP8 operands complicate the sequencing of dependent FP8 WGMMAs.

With respect to attention mechanisms, the imposed layout constraints require specific adaptations to be made in the FP8 algorithm's structure; these are detailed in \cref{sec:algofp8}.

\subsection{Standard Attention and Flash Attention} In line with \citet{dao2022flashattention}, we define \textbf{standard attention} as the attention mechanism executed on the GPU where the intermediary matrices $\mathbf{S}$ and $\mathbf{P}$ are explicitly stored in HBM. The central innovation behind \textsc{FlashAttention} is the utilization of a localized form of the softmax operation, which is employed to circumvent costly memory operations by consolidating the attention computation within a single kernel. The application of local softmax can be seen in lines \ref{code-ws:softmax_start}-\ref{code-ws:softmax_end} of the main loop for the consumer in \cref{alg:flash3_wgmma_ws_only}, along with rescaling blocks of $\mathbf{O}$. For a straightforward demonstration that this process indeed yields $\mathbf{O}$, see \cite[\S 2.3.1]{dao2023flashattention2}.

\section{FlashAttention-3: Algorithm}
\label{sec:algo}

This section presents the \textsc{FlashAttention-3} algorithm, concentrating primarily on the forward computation, while the details for the backward computation are provided in~\cref{sec:algo_ws_bwd}. We begin by outlining the method for incorporating warp-specialization alongside a circular SMEM buffer into the original structure of \textsc{FlashAttention-2}. Next, we discuss how leveraging WGMMA’s asynchronous properties facilitates the creation of a 2-stage GEMM-softmax pipeline with overlapping computation. In addition, we elaborate on the necessary changes for supporting FP8, addressing both the requirements for layout compatibility and the preservation of accuracy through block quantization and non-uniform processing.

\subsection{Producer-Consumer asynchrony through warp-specialization and pingpong scheduling}
\label{sec:algo_ws}

\paragraph{Warp-specialization}
Similar to \textsc{FlashAttention-2}, the forward computation in \textsc{FlashAttention-3} can be highly parallelized along the dimensions of batch size, head count, and query sequence length. Therefore, we focus on presenting the algorithm at the CTA granularity, where each CTA processes a specific tile $\mathbf{Q}_i$ from the query matrix and generates the corresponding output tile $\mathbf{O}_i$. For clarity, the following explanation presents the warp-specialization approach using a circular SMEM buffer, \emph{excluding} any overlapping between GEMM and softmax computations. Consider $d$ as the head dimension and $N$ as the sequence length. Selecting a query block size $B_r$, the query matrix $\mathbf{Q}$ is partitioned into $T_r = \lceil \frac{N}{B_r} \rceil$ segments $\mathbf{Q}_1, .., \mathbf{Q}_{T_r}$.

\begin{algorithm}[H]
    \caption{\small\label{alg:flash3_wgmma_ws_only}\textsc{FlashAttention-3} forward pass \textbf{excluding} intra-consumer overlap -- CTA view}
    \begin{algorithmic}[1]
\REQUIRE Inputs are matrices $\mathbf{Q}_i \in \mathbb{R}^{B_r \times d}$ and $\mathbf{K}, \mathbf{V} \in \mathbb{R}^{N \times d}$ located in HBM, along with the key block size $B_c$, such that $T_c = \lceil \frac{N}{B_c} \rceil$.
\STATE Set up a pipeline structure for synchronizing stages via a circular $s$-stage SMEM buffer and barriers.
\IF {current warpgroup is the producer}
\STATE Free a specified set of registers in advance.
\STATE Begin transferring $\mathbf{Q}_i$ from HBM to shared memory.
\STATE When loading completes, send signal to consumer warpgroup that $\mathbf{Q}_i$ is now in shared memory.
\FOR{$j$ from $0$ to $T_c-1$}
    \STATE Stall until the consumer signals that the $(j\,\%\,s)$ buffer segment is ready for reuse.
    \STATE Initiate data transfer for $\mathbf{K}_j$ and $\mathbf{V}_j$ from HBM to shared memory within the $(j\,\%\,s)$ buffer slot.
    \STATE After successful transfer, inform consumers that $\mathbf{K}_j$ and $\mathbf{V}_j$ are available.
\ENDFOR
\ELSE 
\STATE Allocate previously freed registers based on how many consumer warps are involved.
\STATE Initialize on-chip variables: $\mathbf{O}_i = (0) \in \mathbb{R}^{B_r \times d}$, $\ell_i = 0$, and $m_i = -\infty$ in $\mathbb{R}^{B_r}$.
\STATE Hold until $\mathbf{Q}_i$ is accessible in shared memory.
\FOR{$j$ ranging from $0$ to $T_c-1$}
\STATE Await the presence of $\mathbf{K}_j$ in shared memory.
\STATE Perform the computation $\mathbf{S}_i^{(j)} = \mathbf{Q}_i \mathbf{K}_j^T$ (using SS-GEMM). Commit this operation and wait. \label{alg:ws_only_gemm1}
\STATE Backup $m_i$ as $m_i^{\mathrm{old}}$ and update $m_i = \mathrm{max}(m_i^{\mathrm{old}}, \mathrm{rowmax}(\mathbf{S}_i^{(j)}))$. \label{code-ws:softmax_start}
\STATE Compute the unnormalized attention matrix $\widetilde{\mathbf{P}}_i^{(j)} = \mathrm{exp}(\mathbf{S}_i^{(j)} - m_i)$ and update $\ell_i$ as $\ell_i = \mathrm{exp}(m_i^{\mathrm{old}} - m_i) \ell_i + \mathrm{rowsum}(\widetilde{\mathbf{P}}_i^{(j)})$. \label{code-ws:softmax_end}
\STATE Ensure that $\mathbf{V}_j$ is present in shared memory.
\STATE Update the output: $\mathbf{O}_i = \mathrm{diag}(\mathrm{exp}(m_i^{\mathrm{old}} - m_i))^{-1} \mathbf{O}_i + \widetilde{\mathbf{P}}_i^{(j)} \mathbf{V}_j$ (apply RS-GEMM). Commit and wait for completion. \label{alg:ws_only_gemm2}
\STATE Mark the $(j\,\%\,s)$ buffer region as available for the producer’s next use.
\ENDFOR
\STATE Finalize: normalize output by computing $\mathbf{O}_i = \mathrm{diag}(\ell_i)^{-1} \mathbf{O}_i$ and set $L_i = m_i + \log(\ell_i)$.
\STATE Transfer the finalized block $\mathbf{O}_i$ and value $L_i$ to HBM for their respective positions in $\mathbf{O}$ and $L$.
\ENDIF
\end{algorithmic}
\end{algorithm}

In our implementation of \cref{alg:flash3_wgmma_ws_only} on Hopper, we utilize \verb|setmaxnreg| for allocating and deallocating resources, rely on TMA for loading $\mathbf{Q}_i$ along with $\{\mathbf{K}_j, \mathbf{V}_j\}_{0 \leq j < T_c}$, and apply WGMMA to conduct GEMMs in the primary loop of the consumer. Here, the SS or RS prefixes signify whether the initial operand comes from shared memory or the register file, respectively. When examining the execution of \cref{alg:flash3_wgmma_ws_only}, it is important to recognize that TMA load operations can be dispatched asynchronously and do not block pending the completion of other load requests. In the producer mainloop, the buffer is filled during the initial $s$ iterations, so there are no wait instructions issued during these steps.

\paragraph{Pingpong scheduling} The inherent asynchrony between WGMMA and TMA operations, combined with warp-specialized computation, allows for overlapping the softmax processing in one warpgroup with GEMM computation in another. This is advantageous because, on contemporary hardware accelerators, matmul operations achieve substantially higher throughput than non-matmul functions. For instance, the H100 SXM5 GPU delivers up to 989 TFLOPS for FP16 matmul tasks but is limited to just 3.9 TFLOPS for special functions such as exponential\footnote{According to the CUDA programming guide, each streaming multiprocessor (SM) can execute 16 special function operations per clock cycle. Multiplying 16 by 132 SMs and a clock speed of 1830 MHz results in 3.9 TFLOPS for special functions.} computations required in softmax. When performing the attention forward pass using FP16 with a head dimension of 128, matmul operations generate 512 times more FLOPS relative to exponentials. Nevertheless, exponentials are processed at 256 times lower throughput, meaning they can occupy up to 50\% of the overall cycle time compared to matmul. The gap widens with FP8, as matmul throughput doubles but exponential throughput remains unchanged.

Because the exponential operation is handled by a dedicated hardware component known as the multi-function unit, the optimal scenario would have the exponential computation take place concurrently with the Tensor Cores’ matmul operations. To orchestrate this overlap, we employ synchronization barriers (\texttt{bar.sync} instructions), which enforce an ordering where the GEMMs (specifically, GEMM1—$\mathbf{P} \mathbf{V}$ from one iteration, and GEMM0—$\mathbf{Q} \mathbf{K}^\top$ from the subsequent iteration) belonging to warpgroup 1 are executed prior to those of warpgroup 2. This arrangement ensures that, while warpgroup 2 undertakes its GEMM computations, the softmax operation for warpgroup 1 is processed. Afterwards, the responsibilities switch: warpgroup 1 proceeds with GEMMs as warpgroup 2 executes softmax. This alternating assignment, referred to as “pingpong” scheduling, is depicted in~\cref{fig:pingpong_scheduling}. Although the pingpong execution pattern is not perfectly synchronized in practical implementations, we consistently observe enhanced throughput (e.g., performance rising from 570 TFLOPS to 620–640 TFLOPS for FP16 forward passes with a head dimension of 128 and sequence length of 8192).

\paragraph{Attention variants} In the cases of multi-query attention~\citep{shazeer2019fast} and grouped query attention~\citep{ainslie2023gqa}, we adopt the methodology described in \textsc{FlashAttention-2}, modifying the tensor indexing so as to eliminate redundant copies of $\mathbf{K}$ and $\mathbf{V}$ within HBM.

\subsection{Intra-warpgroup overlapping GEMMs and softmax}

Within a single warpgroup, it is possible to concurrently execute certain softmax instructions alongside selected GEMM instructions. Here, we outline a specific method for achieving such overlap.

In the attention algorithm, operations within the inner loop (main loop) have sequential dependencies that impede parallelization within a single iteration. For example, (local) softmax (lines \ref{code-ws:softmax_start} to \ref{code-ws:softmax_end}) relies on the output $\mathbf{S}_i^{(j)}$ of the first GEMM, while the second GEMM takes its result $\widetilde{\mathbf{P}}_i^{(j)}$ as an operand. Indeed, the wait statements in lines \ref{alg:ws_only_gemm1} and \ref{alg:ws_only_gemm2} of \cref{alg:flash3_wgmma_ws_only} serialize the execution of softmax and GEMMs. However, we can break these dependencies by pipelining across iterations through additional buffers in registers. Pursuing this idea, we propose the following two-stage\footnote{Note that the number of stages of the overlapping scheme is bounded by, but need not equal, the number $s$ of stages in the circular SMEM buffer.} GEMM-softmax pipelining algorithm:

\begin{algorithm}[H]
    \caption{\small\label{alg:flash3_wgmma}\textsc{FlashAttention-3} consumer warpgroup forward pass}
    \begin{algorithmic}[1]
      \REQUIRE Matrices $\mathbf{Q}_i \in \mathbb{R}^{B_r \times d}$, and $\mathbf{K}, \mathbf{V} \in \mathbb{R}^{N \times d}$ located in HBM; key block size $B_c$ specified with $T_c = \lceil \frac{N}{B_c} \rceil$.
      \STATE Adjust register allocation dynamically depending on the number of consumer warps available.
      \STATE Initialize $\mathbf{O}_i = (0) \in \mathbb{R}^{B_r \times d}$, and $\ell_i, m_i = (0), (-\infty) \in \mathbb{R}^{B_r}$ on-chip.
      \STATE Wait until $\mathbf{Q}_i$ and $\mathbf{K}_0$ are in shared memory.
      \STATE Using WGMMA, calculate $\mathbf{S}_{\mathrm{cur}} = \mathbf{Q}_i \mathbf{K}_0^T$. Commit operation, then synchronize.
      \STATE Release the buffer’s initial stage allocated for $\mathbf{K}$.
      \STATE Update $m_{i}$, $\tilde{\mathbf{P}}_{\mathrm{cur}}$, and $\ell_{i}$ from $\mathbf{S}_{\mathrm{cur}}$, and scale $\mathbf{O}_i$ accordingly.
      \FOR{$1 \le j < T_c - 1$}
        \STATE Await loading of $\mathbf{K}_j$ into shared memory.
        \label{code:mainloop_start}
        \STATE Perform WGMMA to compute $\mathbf{S}_{\mathrm{next}} = \mathbf{Q}_i \mathbf{K}_{j}^T$. Commit but do not block.
        \label{code:first_wgmma}
        \STATE Wait for $\mathbf{V}_{j-1}$ to become available in shared memory.
        \STATE Using WGMMA, compute $\mathbf{O}_{i} = \mathbf{O}_{i} + \tilde{\mathbf{P}}_{\mathrm{cur}} \mathbf{V}_{j-1}$. Commit without waiting. 
        \label{code:second_wgmma}
        \STATE Synchronize for completion of WGMMA calculation of $\mathbf{Q}_i \mathbf{K}_{j}^T$.
        \STATE From $\mathbf{S}_{\mathrm{next}}$, update $m_{i}$, $\tilde{\mathbf{P}}_{\mathrm{next}}$, and $\ell_{i}$.
        \label{code:softmax}
        \STATE Wait for WGMMA result of $\tilde{\mathbf{P}}_{\mathrm{cur}} \mathbf{V}_{j-1}$, then apply the scaling to $\mathbf{O}_i$.
        \STATE Release the buffer's $(j\,\%\,s)$th and $(j-1\,\%\,s)$th stages for $\mathbf{K}$ and $\mathbf{V}$ respectively.
        \STATE Copy the contents of $\mathbf{S}_{\mathrm{next}}$ over to $\mathbf{S}_{\mathrm{cur}}$.
        \label{code:mainloop_end}
      \ENDFOR 
      \STATE Wait for loading of $\mathbf{V}_{T_c - 1}$ into shared memory.
      \STATE Execute WGMMA to calculate $\mathbf{O}_{i} = \mathbf{O}_{i} + \tilde{\mathbf{P}}_{\mathrm{last}} \mathbf{V}_{T_c - 1}$. Commit and synchronize at completion.

\STATE Epilogue: Adjust $\mathbf{O}_{i}$ according to $m_{i}$. Determine $L_{i}$ utilizing $m_{i}$ and $\ell_{i}$.
Store $\mathbf{O}_{i}$ and $L_{i}$ in HBM, assigning them as the $i$-th segments of $\mathbf{O}$ and $L$.

\cref{alg:flash3_wgmma} serves to substitute the consumer phase of \cref{alg:flash3_wgmma_ws_only}, thereby forming the complete \textsc{FlashAttention-3} procedure for computations in FP16 format. At a broad level, WGMMA is employed as shorthand for asynchronous GEMM. During the mainloop, demarcated by lines \ref{code:mainloop_start} through \ref{code:mainloop_end}, the second invocation of WGMMA in iteration $j$ (see line \ref{code:second_wgmma}) is concurrently executed alongside the softmax computations associated with iteration $j+1$ (line \ref{code:softmax}).

Although the pipelined architecture presented previously promises improved theoretical performance, several real-world factors must be taken into account:

\paragraph{Compiler reordering} The pseudocode assumes a perfect, deterministic instruction ordering; however, the NVCC compiler often modifies the sequence to optimize for throughput. Such reordering may interfere with the intended interleaving of WGMMA and non-WGMMA instructions, possibly causing anomalous outcomes or undermining performance. Examination of the SASS output confirms that the compiler does, in fact, generate overlapping instructions as anticipated (see Section~\ref{sec:2-stage-sass}).

\paragraph{Register pressure} Achieving maximal performance necessitates minimizing register spilling. The 2-stage pipeline, however, increases the demand for registers since intermediate values and inter-stage context must be retained. In particular, a separate $\mathbf{S}_{\mathrm{next}}$ must reside in registers, resulting in additional usage amounting to $B_r \times B_c \times \text{sizeof}(\text{float})$ per threadblock. This heightened register consumption may compete with optimizations employing larger block sizes—which themselves require more registers. Consequently, optimal parameter choices must be guided by profiling and empirical measurement.

\paragraph{3-stage pipelining} Building upon the 2-stage method above, a further extension to 3 stages is proposed to enable overlap between the second WGMMA and the softmax calculation. Although this variant could boost Tensor Core utilization further, it increases register requirements even more due to the additional pipelining stage. This complicates the balance between tile dimension selection and pipeline depth. For a comprehensive account of the 3-stage approach, including empirical results, refer to~\cref{sec:3-stage}.

\subsection{Low-precision with FP8}
\label{sec:algofp8}

\textbf{Efficiency: layout transformations.} Executing the forward computation of \textsc{FlashAttention-3} using FP8 precision introduces unique difficulties related to layout compatibility, which do not arise when utilizing FP16.

To begin, we observe that input tensors $\mathbf{Q}$, $\mathbf{K}$, and $\mathbf{V}$ are ordinarily arranged in memory with contiguous head dimension. However, the FP8 WGMMA's k-major requirement for the second GEMM necessitates that $\mathbf{V}$—specifically its tiles when placed in SMEM—be contiguous along the sequence length axis. Because the TMA load operation does not alter which dimension is contiguous, achieving this layout demands either (1) transposing $\mathbf{V}$ within GMEM prior to computation, or (2) performing a tile-wise transpose inside the kernel after loading into SMEM. To realize approach (1), we could (1a) integrate the transpose into the epilogue of an earlier stage such as rotary embedding, or (1b) invoke a dedicated transpose kernel as a pre-processing step\footnote{Efficient transpose kernels can run close to the device's maximum bandwidth~\citep{colfax_cutlass_transpose_2024}.}, swapping the strides of the sequence length and head dimensions. Nevertheless, (1a) poses significant challenges for inclusion in standardized libraries, whereas (1b) imposes unacceptable overhead in situations where memory bandwidth dominates, like inference.

Conversely, for FP8 \textsc{FlashAttention-3} we select approach (2). Within the kernel transpose step, we exploit the LDSM (\verb|ldmatrix|) and STSM (\verb|stmatrix|) instructions. These operations allow a warp of threads to cooperatively load data from SMEM into RMEM, and to store data from RMEM to SMEM, each time handling 128-byte segments.\footnote{As noted in the PTX documentation, LDSM and STSM are designed for moving $8 \times 8$ matrices with 16-bit elements \cite[\S 9.7.13.4.15-16]{ptx}. To support FP8, we pack pairs of 8-bit values together, enabling LDSM/STSM to operate under FP8 format. However, the transpose variants of these instructions cannot separate the packed 8-bit elements, which requires some additional register-level manipulation between the LDSM and STSM calls to complete a tiled transpose—these specifics are omitted here.} Employing LDSM/STSM is beneficial in terms of register utilization, as we can schedule their execution within the producer warpgroup, and they support transposed data movement during memory copy. Additionally, it is possible—after the initial iteration—to schedule the transposition of the subsequent $\mathbf{V}$ tile to proceed concurrently with the two WGMMAs processing the previous $\mathbf{V}$ and the current $\mathbf{K}$ tile.

Secondly, we note a distinction between FP16 and FP8 WGMMA in terms of memory layout: the FP32 accumulator in FP8 WGMMA does not align with the memory configuration used for operand A when stored in registers. This difference is illustrated in \cref{fig:rmem_accum} and \cref{fig:rmem_operand}, which show partial layouts where register entries per thread are arranged as indicated. Through the use of byte permutation instructions, the accumulator from the first WGMMA can be reformatted to match the layout expected by a subsequent WGMMA as well as the layout of the $\mathbf{V}$ tile resulting from the in-kernel transpose operation. To be specific, referring to \cref{fig:rmem_accum}, the registers are reordered sequentially as
$$\{ \verb|d0 d1 d4 d5 d2 d3 d6 d7| \},$$
with this rearrangement repeated every 8 bytes across the registers. Logically, this operation permutes the columns of the $\mathbf{P}$ tile (for example, columns $0189$ are now positioned as the initial four columns). To ensure the output tile produced by WGMMA is correct, the in-kernel transpose can be modified to output the $\mathbf{V}$ tile according to a corresponding row permutation.\footnote{Conducting the in-kernel transpose in this manner obviates the need for shuffle instructions to reallocate register ownership among threads, an approach previously discussed in~\citep{colfax_fp8_flashattention_2024}.}

\textbf{Accuracy: block quantization and incoherent processing.} The FP8 (e4m3) representation allocates only 3 bits to the mantissa and 4 bits for the exponent, leading to increased numerical inaccuracies compared to FP16/BF16. Furthermore, in large-scale models, it is common to observe outlier values~\citep{dettmers2208llm,
  sun2024massive} that have much higher magnitudes than the bulk of the parameters, which complicates the quantization process. A standard approach involves per-tensor scaling~\citep{micikevicius2022fp8}, where a single scaling factor is assigned to each tensor (e.g., one each for $\mathbf{Q}$, $\mathbf{K}$, and $\mathbf{V}$).
To address the increased numerical errors in attention arising from FP8, we incorporate two methods:

\begin{enumerate}

\item \textbf{Block quantization}: Here, a single scalar is retained for each block, meaning that tensors $\mathbf{Q}$, $\mathbf{K}$, and $\mathbf{V}$ are partitioned into blocks of dimensions $B_r \times d$ or $B_c \times d$, with each block independently quantized. The quantization process can be integrated with preceding operations in the attention pipeline, such as rotary embedding, without incurring latency penalties (as the rotary embedding operation is primarily constrained by memory bandwidth). Because the \textsc{FlashAttention-3} algorithm inherently performs computations in blocks, it is feasible to rescale each block of $\mathbf{S}$ in accordance with block quantization without adding computational overhead.

\item \textbf{Incoherent processing}: In order to mitigate the influence of outlier values, the matrices $\mathbf{Q}$ and $\mathbf{K}$ are pre-multiplied by a random orthogonal matrix $\mathbf{M}$ prior to FP8 quantization. Owing to the orthogonality of $\mathbf{M}$, i.e., $\mathbf{M} \mathbf{M}^\top = I$, it follows that $(\mathbf{Q}\mathbf{M})(\mathbf{K}\mathbf{M})^\top = \mathbf{Q}\mathbf{K}^\top$, demonstrating that such transformation preserves the integrity of the attention computation. This approach distributes outlier values more evenly because each element of $\mathbf{Q}\mathbf{M}$ or $\mathbf{K}\mathbf{M}$ is constructed as a random linear combination of entries from $\mathbf{Q}$ or $\mathbf{K}$, thereby decreasing quantization-induced inaccuracies. Consistent with \citet{chee2024quip} and~\citet{tseng2024quip}, $\mathbf{M}$ is selected as the product of random diagonal matrices containing $\pm 1$ and a Hadamard matrix, which enables matrix multiplication in $O(d \log d)$ rather than $O(d^2)$, and can also be amalgamated with rotary embedding operations without extra computational expense.
\end{enumerate}

Empirical results show that these strategies can lower numerical errors by as much as a factor of 2.6, as demonstrated in \cref{sec:numerical_error}.

\section{Empirical Validation}
\label{sec:exp}

To construct \textsc{FlashAttention-3} and assess its performance in terms of speed and correctness, we employ primitives provided by CUTLASS~\citep{Thakkar_CUTLASS_2023}, including the WGMMA and TMA abstractions.

\begin{itemize}

\item \textbf{Benchmarking attention.} We evaluate the runtime performance of \textsc{FlashAttention-3} for a variety of sequence lengths and directly contrast these results with a baseline PyTorch implementation, as well as with \textsc{FlashAttention-2}, \textsc{FlashAttention-2} in Triton (leveraging H100-optimized instructions), and the cuDNN vendor-tuned version of \textsc{FlashAttention-2} for H100 GPUs. Our experiments show that \textsc{FlashAttention-3} achieves speeds up to 2.0$\times$ greater than \textsc{FlashAttention-2}, and improves upon the Triton variant by 1.5$\times$. Additionally, \textsc{FlashAttention-3} attains up to 740 TFLOPs/s, which corresponds to 75\% of the H100 GPU’s peak theoretical throughput.
\item \textbf{Ablation study.} Through targeted ablations, we demonstrate that introducing warp-specialization and GEMM-softmax pipelining provides significant speed benefits in \textsc{FlashAttention-3}.
\item \textbf{Accuracy of FP8 attention.} Our results show that block quantization and incoherent processing strategies decrease the numerical error in FP8 \textsc{FlashAttention-3} by a factor of 2.6$\times$.
\end{itemize}

\subsection{Benchmarking Attention}
\label{subsec:benchmark_attn}

We evaluate the execution time of various attention mechanisms on an H100 80GB SXM5 GPU across multiple scenarios (using or not using a causal mask, and with head dimensions set to either 64 or 128), processing FP16 input data. The findings are presented in~\cref{fig:benchmark_attn_fwd} and~\cref{fig:benchmark_attn_bwd}. These results indicate that \textsc{FlashAttention-3} achieves approximately 1.5-2.0$\times$ speedup over \textsc{FlashAttention-2} during the forward computation, and 1.5-1.75$\times$ during the backward computation. Relative to a conventional attention implementation, \textsc{FlashAttention-3} reaches accelerations ranging from 3$\times$ to 16$\times$. Notably, for sequence lengths of 1,000 tokens and greater, \textsc{FlashAttention-3} performs faster than even the vendor-provided library (cuDNN—closed source) tailored for H100 GPUs.

\paragraph{Benchmark settings:} We vary the sequence length as 512, 1k, ..., 16k, and set batch size so that the total number of tokens is 16k. We set the hidden dimension to 2048, and head dimension to be either 64, 128, or 256 (i.e., 32 heads, 16 heads, or 8 heads). To calculate the FLOPs of the forward pass, we use:

\begin{equation*}
  4 \cdot \text{seqlen}^2 \cdot \text{head dimension} \cdot \text{number of heads}.
\end{equation*}

By implementing causal masking, the total is halved since roughly 50% of the computations are performed. For the backward pass, we estimate the FLOPs by scaling the forward pass count by 2.5, reflecting that the backward step involves five matrix multiplications compared to two in the forward step, largely as a result of recomputation.

Additionally, we assess the forward pass runtime for FP8 within comparable experimental conditions. The outcomes corresponding to headdim 256 are presented in ~\cref{fig:benchmark_attn_fp8}, with comprehensive data provided in \cref{sec:benchmark_attn_fp8_full}.

\subsection{Ablation Study: 2-Stage Pipelining Experiments}
\label{sec:2-stage-pipelining-experiments}

We investigate the effects of removing both the two-stage WGMMA-softmax pipeline and the warp-specialization technique on non-causal FP16 \textsc{FlashAttention-3}, using constant settings $\{\text{batch}, \text{seqlen}, \text{nheads}, \text{hdim}\} = \{ 4, 8448, 16, 128\}$. The findings, displayed in~\cref{table:ablation_pipelining}, demonstrate that our algorithmic enhancements—which combine asynchronous execution via warp-specialization and concurrent GEMM-softmax computation—produce a substantial performance gain, increasing throughput from 570 to 661 TFLOPs.

\subsection{Numerical Error Validation}
\label{sec:numerical_error}

Given the growing attention to the numerical error associated with \textsc{FlashAttention}~\citep{golden2024flash}, we conduct a comparative analysis involving \textsc{FlashAttention-2}, \textsc{FlashAttention-3}, and a conventional attention implementation, using a reference version in FP64 as a baseline. In order to replicate the presence of outlier features and activations observed in LLMs~\citep{dettmers2208llm, sun2024massive}, we generate the elements of $\mathbf{Q}, \mathbf{K}, \mathbf{V}$ according to the following statistical distribution:

\begin{equation*}
  \mathcal{N}(0, 1) + \mathcal{N}(0, 100) \cdot \mathrm{Bernoulli}(0.001).
\end{equation*}

In this setup, the majority of entries are sampled from a normal distribution with mean zero and standard deviation 1, but for 0.1\% of entries, we introduce an additional independent component drawn from a normal distribution with standard deviation 10. We report the root mean squared error (RMSE) results in~\cref{table:numerical_error}. When operating in FP16, both \textsc{FlashAttention-2} and \textsc{FlashAttention-3} produce RMSE values that are 1.7$\times$ lower than those from the conventional implementation, benefiting from the retention of intermediate (softmax) values in FP32. The baseline FP8 attention employs per-tensor scaling, with the matrix multiplication accumulator in FP32 and the intermediate softmax outputs in FP16. Due to the utilization of block quantization and incoherent processing, \textsc{FlashAttention-3} in FP8 attains an accuracy improvement of 2.6$\times$ over this baseline.

\section{Dicussion, Limitations, Conclusion}
\label{sec:discussion}

Through the development of \textsc{FlashAttention-3}, we illustrate that advancements in programming paradigms and leveraging hardware innovations like asynchrony and low-precision computation substantially enhance both performance and reliability in attention mechanisms. The attention operation achieves a 1.5-2.0$\times$ acceleration over \textsc{FlashAttention-2} and offers a 2.6$\times$ reduction in FP8 numerical error relative to typical per-tensor quantization methods. There remain aspects to improve, including optimization for LLM inference workloads, adopting a persistent kernel architecture for the FP8 kernel,\footnote{In our benchmarking results, FP16 \textsc{FlashAttention-3} incorporates both a persistent kernel and a load balancing scheme; however, FP8 \textsc{FlashAttention-3} currently lacks these features. This omission contributes to FP8 \textsc{FlashAttention-3}'s inferior performance on short sequences and causal masking tasks versus FP8 cuDNN kernels.} and further investigation into the implications of low-precision attention during large-scale model training. While our experiments centered on Hopper GPUs, we anticipate the methodologies presented will generalize to other accelerators. Our goal is that improvements in speed and accuracy of fundamental operations such as attention will enable new possibilities for applications involving extended contexts.

\end{document}